% swarm_blocker6_v2_agent14_delta5_reconstruction.tex
% Agent 14 / Blocker 6 v2 --- conditional Delta_5 compatibility for the
% perturbative E_3 factorisation algebra of 6d hCS on K3 x E.
% Standalone preamble; no other-file modification.
% Cross-refs (to manuscript, by name only): thm:6d-hcs-factorisation-algebra-k3xe,
% thm:g-delta5-is-sp-k3, subsec:k3xE-delta5-specialisation,
% chapters/examples/k3e_bkm_chapter.tex.

\documentclass[11pt]{amsart}

\usepackage{amsmath,amssymb,amsthm,mathrsfs}
\usepackage[margin=1in]{geometry}
\usepackage{hyperref}
\usepackage{microtype}

\providecommand{\C}{\mathbb{C}}
\providecommand{\R}{\mathbb{R}}
\providecommand{\Z}{\mathbb{Z}}
\providecommand{\N}{\mathbb{N}}
\providecommand{\Q}{\mathbb{Q}}
\providecommand{\HH}{\mathrm{HH}}
\providecommand{\Perf}{\mathrm{Perf}}
\providecommand{\Coh}{\mathrm{Coh}}
\providecommand{\cE}{\mathcal{E}}
\providecommand{\cF}{\mathcal{F}}
\providecommand{\cO}{\mathcal{O}}
\providecommand{\cA}{\mathcal{A}}
\providecommand{\cZ}{\mathcal{Z}}
\providecommand{\cC}{\mathcal{C}}
\providecommand{\cH}{\mathcal{H}}
\providecommand{\cL}{\mathcal{L}}
\providecommand{\cM}{\mathcal{M}}
\providecommand{\cP}{\mathcal{P}}
\providecommand{\cD}{\mathcal{D}}
\providecommand{\cR}{\mathcal{R}}
\providecommand{\cS}{\mathcal{S}}
\providecommand{\cT}{\mathcal{T}}
\providecommand{\dbar}{\bar{\partial}}
\providecommand{\frakg}{\mathfrak{g}}
\providecommand{\frakn}{\mathfrak{n}}
\providecommand{\frakh}{\mathfrak{h}}
\providecommand{\fsl}{\mathfrak{sl}}
\providecommand{\fu}{\mathfrak{u}}
\providecommand{\Sp}{\mathrm{Sp}}
\providecommand{\SO}{\mathrm{SO}}
\providecommand{\tr}{\mathrm{tr}}
\providecommand{\Tr}{\mathrm{Tr}}
\providecommand{\Obs}{\mathrm{Obs}}
\providecommand{\Sym}{\mathrm{Sym}}
\providecommand{\End}{\mathrm{End}}
\providecommand{\Hom}{\mathrm{Hom}}
\providecommand{\id}{\mathrm{id}}
\providecommand{\hCS}{\mathrm{hCS}}
\providecommand{\BV}{\mathrm{BV}}
\providecommand{\Fact}{\mathrm{Fact}}
\providecommand{\Eone}{E_{1}}
\providecommand{\Etwo}{E_{2}}
\providecommand{\Ethree}{E_{3}}
\providecommand{\Ran}{\mathrm{Ran}}
\providecommand{\Heis}{\mathsf{Heis}}
\providecommand{\Muk}{\Lambda_{\mathrm{Muk}}}
\providecommand{\Gdelta}{\mathfrak{g}_{\Delta_{5}}}
\providecommand{\PhiFA}{\Phi^{\mathrm{FA}}}
\providecommand{\SpCh}{\mathrm{Sp}^{\mathrm{ch}}}
\providecommand{\kch}{\kappa_{\mathrm{ch}}}
\providecommand{\kcat}{\kappa_{\mathrm{cat}}}
\providecommand{\kBKM}{\kappa_{\mathrm{BKM}}}
\providecommand{\kfib}{\kappa_{\mathrm{fiber}}}
\providecommand{\Im}{\operatorname{Im}}
\providecommand{\Re}{\operatorname{Re}}
\providecommand{\Vol}{\operatorname{Vol}}
\providecommand{\mult}{\operatorname{mult}}
\providecommand{\BorLift}{\mathrm{BorLift}}
\providecommand{\Grit}{\mathrm{Grit}}
\providecommand{\ClaimStatusLemma}{\textsc{[lemma]}}
\providecommand{\ClaimStatusCorollary}{\textsc{[corollary]}}
\providecommand{\ClaimStatusConjectured}{\textsc{[conjectural]}}
\providecommand{\ClaimStatusHeuristic}{\textsc{[heuristic]}}
\providecommand{\ClaimStatusDerivation}{\textsc{[first-principles]}}
\providecommand{\ClaimStatusDefinition}{\textsc{[definition]}}
\providecommand{\ClaimStatusOpen}{\textsc{[open]}}
\providecommand{\ClaimStatusConditional}{\textsc{[conditional]}}

\theoremstyle{plain}
\newtheorem{theorem}{Theorem}[section]
\newtheorem{proposition}[theorem]{Proposition}
\newtheorem{lemma}[theorem]{Lemma}
\newtheorem{corollary}[theorem]{Corollary}
\newtheorem{conjecture}[theorem]{Conjecture}
\theoremstyle{definition}
\newtheorem{definition}[theorem]{Definition}
\newtheorem{construction}[theorem]{Construction}
\theoremstyle{remark}
\newtheorem{remark}[theorem]{Remark}
\newtheorem{attack}[theorem]{Attack}
\newtheorem{heal}[theorem]{Heal}

\title{Conditional $\Delta_{5}$ compatibility for the perturbative $E_{3}$-FA on
$K3\times E$}
\author{Raeez Lorgat}
\date{2026-04-27}

\begin{document}
\maketitle

\begin{abstract}
The Igusa--Borcherds shadow $\Delta_{5}(Z)$ at $K3\times E$ supplies the
weight $\kBKM(\Delta_{5})=c_{\phi_{0,1}}(0,0)/2=10/2=5$ via the
Gritsenko--Nikulin denominator identity. Independently, the perturbative
BV theory of $6$d holomorphic Chern--Simons on $K3\times E$ assembles
into an $\Ethree$-factorisation algebra
$\Obs^{q,\hCS}(K3\times E,\frakg)$ on the K\"ahler product, locally $E_3$
on polydisks and $E_3$-formal up to the $\mathrm{GRT}_{1}(\Q)$-torsor of
associators. Specialising the Stage-$1$ output along the K3 fibre and
restricting to the elliptic base gives a candidate $\Eone$-chiral
algebra $\cA_{E}^{\hCS}$ on $E$. This note does not reconstruct
$\Delta_{5}$ from perturbative hCS. It records the conditional
comparison needed for the graded vacuum character to match the
meromorphic reciprocal denominator, the extra Hall/vanishing-cycle data
needed to read imaginary-root multiplicities, and the external
automorphic input supplied by Borcherds--Gritsenko or
Oberdieck--Pixton. The reciprocal $\Delta_5^{-1}$ has negative-weight
meromorphic Fourier--Jacobi coefficients; they are not positive-weight
cusp coefficients.
\end{abstract}

\section{Setup and the question}
\label{sec:setup}

\textsc{Scope.} $X = K3\times E$ closed K\"ahler Calabi--Yau threefold
with the product Yau metric and elliptic factor
$E = \C/(\Z+\tau\Z)$, $\Im\tau>0$. The gauge dg-Lie algebra $\frakg$ is
a simple complex Lie algebra with vanishing cubic Casimir
$d^{abc}_{\frakg}=0$ (as in $\fsl_{2}$ or $E_{8}$); the cubic-anomaly
obstruction of higher-rank cases is irrelevant to the present
compatibility problem. The Stage-$1$ output is the $\hbar$-adic
$\Ethree$-factorisation algebra
\[
\Obs^{q,\hCS}(K3\times E,\frakg) \;\in\;
\mathrm{FA}^{\mathrm{hol}}(K3\times E;\mathrm{Ch}^{\hbar})
\]
of perturbative BV observables of the $6$d holomorphic Chern--Simons
action $S_{\hCS}=\int_{X}\Omega_{X}^{(3,0)}\wedge\tr(A\,\dbar A
+\tfrac{2}{3}A\wedge A\wedge A)$ at one-loop, with the wave-function
counter-term $-C_{2}(\frakg)/(2\pi)^{3}$ absorbed into the
holomorphic measure (manuscript Theorem
\texttt{thm:6d-hcs-factorisation-algebra-k3xe}, item (i)). $\Sigma_{2}$
denotes the K3 fibre, $C=E$ the elliptic base.

\textsc{Question.} The Igusa cusp form
$\Phi_{10}(Z)=\Delta_{5}(Z)^{2}$ on $\Sp_{4}(\Z)$ controls the
Oberdieck--Pixton DT partition function
$Z^{K3\times E}_{\mathrm{DT}}=C\cdot\Delta_{5}(Z)^{-2}$
\cite{OberdieckPixton2019}; the Gritsenko--Nikulin lift
$\Delta_{5}=\Grit(\eta^{9}\vartheta_{1})=\BorLift(\phi_{0,1})$
\cite{GritsenkoNikulin1995,Borcherds1998} is a Siegel modular form of
weight $5$ on $\Sp_{4}(\Z)$ with multiplier $v_{\Delta_{5}}$, and
Lorgat 2020 \cite{Lorgat2020} expresses it as an explicit theta-product
$\Delta_{5}=\prod_{(a,b)}v_{a,b}$ over a tower of theta-block coordinates,
with $f(1,1,1)=64$ leading Fourier coefficient.

The question of this note: starting only from the $\hbar$-adic
$\Ethree$-FA $\Obs^{q,\hCS}(X,\frakg)$, can one reconstruct
$\Delta_{5}$? The healed answer is no. One obtains, at most,
conditional compatibility with the reciprocal denominator after the
hCS--Hall/Borcherds comparison and the external automorphic completion
are supplied.

\subsection{The two-stage shadow}
\label{subsec:two-stage-shadow}

The two-stage factorisation
$\Phi_{3}=\SpCh_{\Sigma_{2},C}\circ\PhiFA_{3}$ realises the Igusa--%
Borcherds image of $K3\times E$ as the composite of a canonical $E_{3}$-FA
quantisation and an $\Eone$-chiral specialisation. Stage-$1$ supplies
$\Obs^{q,\hCS}$ at the level of explicit chain-level data; Stage-$2$ at
$(\Sigma_{2},C)=(K3,E)$ is the factorisation-homology integration
$\int_{K3}(-)$ along the K3 fibre, restricted to the elliptic base.

\begin{definition}[The chiral algebra on $E$]
\label{def:chiral-algebra-on-E}
\ClaimStatusDefinition{}
With $X = K3\times E$, projection $p_{E}\colon X\to E$, and the
\v{C}ech--Weiss descent of
$\Obs^{q,\hCS}(X,\frakg)$ on the polydisk atlas of $X$, the
\emph{chiral algebra on $E$} is
\[
\cA_{E}^{\hCS}
\;:=\;
(p_{E})_{\ast}\Obs^{q,\hCS}(X,\frakg)
\;\simeq\;
\bigl(\textstyle\int_{K3}\Obs^{q,\hCS}\bigr)_{E},
\]
the cosheaf pushforward to $E$ of the $\Ethree$-FA observables along
the K3 fibre. The factorisation--homology integration
$\int_{K3}(-)$ is the Lurie symmetric-monoidal functor
$\mathrm{Alg}_{\Ethree}\to\mathrm{Alg}_{\Eone}$
restricted to the K3 fibre.
\end{definition}

\textsc{Test.} Compute the graded vacuum character
$\Tr_{\mathrm{vac}}q^{L_{0}-c/24}$ of $\cA_{E}^{\hCS}$ as a formal
$q$-series; compare, after supplying the Hall/Borcherds comparison, with
the Fourier--Jacobi expansion of $\Delta_{5}(2Z)^{-1}$ along
$z\to 0$, that is, with the elliptic
specialisation
\[
\Delta_{5}(2Z)^{-1}\big|_{E}
\;=\;
\sum_{m\geq -1}\psi_{m,-5}^{\mathrm{mer}}(\tau,z)\,p^{m},
\qquad p:=e^{2\pi i\sigma},
\]
where $\sigma$ is the second period and
$\psi_{m,-5}^{\mathrm{mer}}$ are meromorphic Jacobi coefficients of
reciprocal weight $-5$ in the chosen Klingen chamber.

\section{Fourier--Jacobi expansion of $\Delta_{5}$}
\label{sec:fj-expansion-delta5}

We record the Fourier--Jacobi side of the comparison. The Borcherds
product representation \cite{Borcherds1998}, in the Gritsenko--Nikulin
normalisation \cite{GritsenkoNikulin1995},
\begin{equation}
\label{eq:borcherds-product-delta5}
\Delta_{5}(Z)
=
e^{2\pi i(\rho,Z)}\prod_{\alpha>0}
\bigl(1-e^{2\pi i(\alpha,Z)}\bigr)^{\mult(\alpha)},
\quad
\mult(\alpha)
=
c_{\phi_{0,1}}\!\bigl(-\tfrac{1}{2}(\alpha,\alpha),\ell_{\alpha}\bigr),
\end{equation}
with Weyl vector $\rho=\tfrac{1}{2}(\delta_{1}+\delta_{2}+\delta_{3})$
of squared norm $(\rho,\rho)=-3/2$, three real even simple roots
$\delta_{1},\delta_{2},\delta_{3}$, and imaginary-root multiplicities
read from the weak Jacobi form
\[
\phi_{0,1}(\tau,z)
=
4\sum_{\sigma\in\{2,3,4\}}\bigl(\theta_{\sigma}(\tau,z)/\theta_{\sigma}(\tau,0)\bigr)^{2},
\qquad
c_{\phi_{0,1}}(0,0)=10.
\]
Hence $\kBKM(\Delta_{5})=c_{\phi_{0,1}}(0,0)/2=5$.

\begin{lemma}[Fourier--Jacobi expansion]
\label{lem:fj-expansion-delta5-inverse}
\ClaimStatusLemma{}
The reciprocal $\Delta_{5}(2Z)^{-1}$ admits a Fourier--Jacobi expansion
along $E$ with polar leading term
$p^{-1}(\eta^{9}\vartheta_{1})^{-1}$ and successive meromorphic
Jacobi coefficients of reciprocal weight $-5$ in the Klingen chamber.
These coefficients are not elements of $J_{5,m}^{\mathrm{cusp}}$.
They become BKM character data only after the denominator identity and
the Hall/Borcherds comparison are supplied.
\end{lemma}

\begin{proof}
The Fourier--Jacobi expansion of a Siegel modular form
$F$ of weight $k$ on $\Sp_{4}(\Z)$ (with respect to the Klingen
parabolic) is $F(Z)=\sum_{m}\phi_{m,k}(\tau,z)p^{m}$ with
$\phi_{m,k}\in J_{k,m}$. For $\Delta_{5}$ the leading term is
$\eta^{9}\vartheta_{1}\cdot p$, whence $\Delta_{5}(2Z)^{-1}$ has leading
$p^{-1}(\eta^{9}\vartheta_{1})^{-1}$ and transforms with weight $-5$
and inverse multiplier. The denominator identity
\eqref{eq:borcherds-product-delta5} identifies the reciprocal
denominator with a Weyl--Kac--Borcherds character after the BKM datum is
already known; it does not turn the reciprocal coefficients into
positive-weight cusp forms
\cite{Borcherds1998,GritsenkoNikulin1995}.
\end{proof}

\section{The graded character of $\cA_{E}^{\hCS}$}
\label{sec:graded-char-AE-hcs}

The pushforward to $E$ of $\Obs^{q,\hCS}(K3\times E,\frakg)$ is
computable in BV perturbation theory. The harmonic zero-mode
contribution decomposes as
\[
\cZ_{X}^{\frakg}
\;\cong\;
\frakg^{\oplus h^{0,0}+h^{0,1}+h^{0,2}+h^{0,3}}
\;=\;
\frakg^{4},
\]
since on $K3\times E$ the Hodge numbers $(h^{0,q})_{q=0}^{3}=(1,1,1,1)$
by K\"unneth on $h^{\bullet,0}(K3)=(1,0,1)$ and $h^{\bullet,0}(E)=(1,1)$.
Each higher-order propagator pairing is the Costello--Li renormalised
$\dbar$-Green's kernel on $X$, with infrared support concentrated on the
diagonal of the K3 fibre after $\int_{K3}$.

\begin{proposition}[Vacuum character of $\cA_{E}^{\hCS}$ to one loop]
\label{prop:vacuum-char-AE-one-loop}
\ClaimStatusConditional{}\textup{(}admissible $\frakg$, perturbative,
formal in $\hbar$, item--scoped to one-loop level\textup{)}\\
After a gauge-compatible pushforward to $E$ and a fixed
normal-ordering, the one-loop graded vacuum character of
$\cA_{E}^{\hCS}$ is expected to have a free-field product shape
\[
\Tr_{\mathrm{vac}}^{\cA_{E}^{\hCS}}\bigl(q^{L_{0}-c/24}\bigr)
\;=\;
q^{-c/24}\prod_{n\geq 1}(1-q^{n})^{-a_{n}}
\;+\;O(\hbar^{2}),
\]
where the integers $a_n$ must be computed from the pushed-forward
Dolbeault modes and their renormalised one-loop contractions. They are
not, without the Hall/vanishing-cycle comparison, imaginary-root
multiplicities of $\Gdelta$. The central charge has the expected
Heisenberg--Cartan leading term
\[
c \;=\; 3\cdot\dim\frakg\,\cdot \chi(K3)/24 \;+\; O(\hbar)
\;=\;
3\cdot\dim\frakg \;+\; O(\hbar).
\]
\end{proposition}

\begin{proof}[Sketch]
The pushforward $(p_{E})_{\ast}\Obs^{q,\hCS}$ on $E$ is governed by
\cite{CostelloGwilliam2017} Chapter 5 after a compact pushforward and
renormalisation package are fixed. At zero loops, zero-modes contribute
the $\frakg^{4}$-Gaussian factor; the
$\frakg$-Cartan piece survives the K\"unneth pairing of
$(h^{0,0},h^{0,1})$ across $(K3,E)$ and produces a chiral free boson
on $E$ in each $\frakg$-direction. At one loop, the remaining mode
determinants define the exponents $a_n$ after normal-ordering. This
proves only the expected product shape under the stated hypotheses; it
does not identify the $a_n$ with BKM root multiplicities.
\end{proof}

\subsection{Comparison with $\Delta_{5}(2Z)^{-1}|_{E}$}
\label{subsec:compare-delta5-side}

\begin{theorem}[Conditional reciprocal-denominator comparison target]
\label{thm:perturbative-match-delta5}
\ClaimStatusConditional{}\textup{(}principal locus plus
hCS--Hall/Borcherds trace comparison\textup{)}\\
On the principal $K3\times E$ locus, after supplying the BL-2/4/6
witnesses, the oriented hCS--Hall comparison, and the module-level
trace comparison, the graded vacuum character of $\cA_{E}^{\hCS}$ has
the reciprocal-denominator target
\[
\Tr_{\mathrm{vac}}^{\cA_{E}^{\hCS}}\bigl(q^{L_{0}-c/24}\bigr)
\;\leadsto\;
\Delta_{5}(2Z)^{-1}\big|_{E}
\quad\text{in }\Q[\![q,\hbar]\!].
\]
This note does not prove equality to all perturbative orders.
\end{theorem}

\begin{proof}[Sketch]
The chiral side of the equality is the manuscript theorem
\texttt{thm:g-delta5-is-sp-k3}: the Stage-$2$ specialisation
$\SpCh_{K3,E}(\PhiFA_{3}(\Perf(K3\times E)))$ is the vertex envelope
$U_{\mathrm{ch}}(\Gdelta)$, whose vacuum character is the reciprocal
denominator by the Weyl--Kac--Borcherds character formula. The
factorisation-algebra side $\cA_{E}^{\hCS}$ reaches that chiral algebra
only after the hCS-to-Hall comparison and trace compatibility are
constructed. Proposition~\ref{prop:vacuum-char-AE-one-loop} is at most
a one-loop shape witness. The $\mathrm{GRT}_{1}(\Q)$-torsor of
associators does not by itself force equality of all higher-loop
characters.
\end{proof}

\section{External automorphic completion via Borcherds--Weyl rigidity}
\label{sec:reconstruction-borcherds-weyl}

The comparison target of Theorem~\ref{thm:perturbative-match-delta5}
does not extract imaginary-root multiplicities from hCS by itself.
Those multiplicities are BKM/Hall data. A factorisation-algebra
dimension count becomes a multiplicity statement only after the
Hall/vanishing-cycle comparison identifies the graded hCS trace with the
Borcherds denominator trace.

\begin{lemma}[Imaginary-root multiplicities from FA characters]
\label{lem:imag-root-mult-from-FA}
\ClaimStatusConditional{}
For each Heegner divisor $H_{D}\subset\Sp_{4}(\Z)\backslash\mathbb{H}_{2}$
of discriminant $D=4N-\ell^{2}$, suppose the oriented hCS--Hall
comparison, vanishing-cycle normalisation, stable-envelope transport,
and Hall--Borcherds denominator comparison have been supplied. Then the
multiplicity in \eqref{eq:borcherds-product-delta5} is compared with
the corresponding graded hCS trace by
\[
\mult(\alpha_{D})
\;=\;
\dim_{\C}\bigl(\cA_{E}^{\hCS}\bigr)^{(D)}_{q^{N},\zeta^{\ell}},
\]
where $(\cA_{E}^{\hCS})^{(D)}$ is the Heegner-grading-$D$ subspace of
the FA observables under the Maulik--Okounkov stable-envelope
filtration on $H^{\bullet}_{T}(\cM_{\Coh(K3\times E)})$ pushed to $E$,
and the $(q^{N},\zeta^{\ell})$ refers to the Fourier--Jacobi bigrading.
Without these comparison data, the right-hand side is only a graded
dimension and not an imaginary-root multiplicity.
\end{lemma}

\begin{proof}[Sketch]
By the Borcherds singular-theta lift, the discriminant-$D$ Heegner
divisor $H_{D}$ has Chern class
$c_{1}(\cL^{\Delta_{5}})|_{H_{D}}=\sum_{N,\ell}c_{\phi_{0,1}}(N,\ell)\,
[H_{4N-\ell^{2}}]$ on the Siegel modular threefold (Bruinier 2002
\cite{Bruinier2002} Theorem 5.2, Kudla--Millson 1990 \cite{Kudla1997}
modularity-of-generating-series). Identifying these Chern classes with
characters of FA observables on $E$ is precisely the extra
hCS--Hall/Borcherds obligation. Bruinier's Heegner-Chern-class
reciprocity \cite{Bruinier2002} Proposition 5.1 pins the automorphic
side; it does not construct the hCS comparison.
\end{proof}

\begin{theorem}[Borcherds--Weyl rigidity after external multiplicities]
\label{thm:borcherds-weyl-rigidity-reconstruction}
\ClaimStatusConditional{}\textup{(}principal locus, Borcherds 1998
Theorem~13.3 hypotheses, external multiplicity input\textup{)}\\
If the imaginary-root multiplicities are supplied by the
Hall/Borcherds comparison of Lemma~\ref{lem:imag-root-mult-from-FA},
then Borcherds--Weyl rigidity determines the automorphic denominator
as the weight-$5$ Siegel modular form $\Delta_{5}(Z)$ with multiplier
$v_{\Delta_{5}}$. The rigidity step is automorphic; it is not a
reconstruction from perturbative hCS alone.
\end{theorem}

\begin{proof}[Sketch]
By Borcherds 1998 \cite{Borcherds1998} Theorem 13.3, a generalised
Kac--Moody superalgebra is determined up to isomorphism by its even
real simple roots, its imaginary-root multiplicities (read off the
denominator), and its grading lattice. For the present case, the lattice
is $\Lambda^{2,1}_{II}\subset\Lambda^{3,2}$ (Gritsenko--Nikulin Weyl
polygon $\cP_{II}$ \cite{GritsenkoNikulin1995}), the three real simple
roots are $\delta_{1},\delta_{2},\delta_{3}$ with Weyl vector
$\rho=\tfrac{1}{2}(\delta_{1}+\delta_{2}+\delta_{3})$, and the
imaginary-root multiplicities are supplied by the external comparison
hypotheses. The denominator identity
\eqref{eq:borcherds-product-delta5} then gives $\Delta_{5}(Z)$ as a
Borcherds product. The rigidity of the singular-theta lift on
the singular weight $\kBKM=c_{\phi_{0,1}}(0,0)/2=5$ within
$M_{5}(\Sp_{4}(\Z),v_{\Delta_{5}})$, a one-dimensional space spanned by
$\Delta_{5}$ alone (Gritsenko 1999 \cite{GritsenkoNikulin1995}, Lorgat
2020 \cite{Lorgat2020}), forces the unique automorphic completion.
\end{proof}

\section{Attack--heal cycles}
\label{sec:attack-heal-cycles}

We isolate three obstructions and resolve each. Each cycle takes the
strongest counterexample as the attack, then heals the proof with the
appropriate ambient hypothesis.

\subsection{Cycle 1: divergence of the Fourier--Jacobi sum}

\begin{attack}
\label{attack:fj-divergence}
The Fourier--Jacobi expansion $\sum_{m}\phi_{m,5}p^{m}$ converges only
in the Klingen parabolic chamber $\Im\sigma > C(\Im\tau,\Re z)$ for
some explicit threshold $C$. The graded character of $\cA_{E}^{\hCS}$
is by construction a $q$-series on $E$ alone; its $p$-expansion is at
best formal. Moreover for the reciprocal the coefficients are
meromorphic negative-weight Jacobi coefficients, not
$\phi_{m,5}\in J_{5,m}^{\mathrm{cusp}}$. Thus the comparison target of
Theorem~\ref{thm:perturbative-match-delta5} cannot by itself pin a
unique $\Delta_{5}$.
\end{attack}

\begin{heal}
\label{heal:fj-divergence}
Borcherds--Weyl rigidity \cite{Borcherds1998} Theorem 13.3 applies
after the imaginary-root multiplicities
$\{\mult(\alpha_D)\}_{D\geq 0}$ are known as BKM/Hall data.
Lemma~\ref{lem:imag-root-mult-from-FA} is therefore conditional: a finite
FA subspace dimension becomes a multiplicity only after the
hCS--Hall/Borcherds comparison identifies the trace. Once those
multiplicities are externally supplied, the singular-theta lift gives
the automorphic completion.
\end{heal}

\subsection{Cycle 2: gauge dependence of the norm-ordered character}

\begin{attack}
\label{attack:gauge-dependence-character}
The graded character of $\cA_{E}^{\hCS}$ depends on the choice of
Costello--Li propagator within the $\mathrm{GRT}_{1}(\Q)$-torsor of
associators (manuscript \texttt{thm:6d-hcs-factorisation-algebra-k3xe}).
Different propagators produce different normal-ordering of the
factorisation observables, hence different graded characters. The
match in Theorem~\ref{thm:perturbative-match-delta5} therefore depends
on the choice of associator, and is not intrinsic.
\end{attack}

\begin{heal}
\label{heal:gauge-dependence-character}
The heal is to keep gauge invariance as a required datum. A
$\mathrm{GRT}_{1}(\Q)$ gauge transformation may preserve cohomological
traces only after the $L_0$-grading, completion, and normal-ordering
are transported compatibly. The character comparison must therefore be
stated modulo a chosen gauge-compatible hCS--Hall trace package. The
number $\kBKM=c_{\phi_{0,1}}(0,0)/2=5$ is pinned by the external
Borcherds input, not by the propagator torsor.
\end{heal}

\subsection{Cycle 3: singular-theta vs.\ shuffle-coproduct mismatch}

\begin{attack}
\label{attack:singular-theta-shuffle-mismatch}
The Borcherds singular-theta lift is a contour integral of a vector-%
valued theta series against a Maass--Borel--Eisenstein kernel
\cite{Borcherds1998}. The Schiffmann--Vasserot shuffle coproduct
\cite{SchiffmannVasserot2013} on $\mathrm{CoHA}_{K3\times E}$ is, by
contrast, the extension-pullback on the Hall stack. These are
structurally different operations. Lemma~\ref{lem:imag-root-mult-from-FA}
identifies them at the level of Heegner-divisor evaluation; this
identification cannot be unconditional.
\end{attack}

\begin{heal}
\label{heal:singular-theta-shuffle-mismatch}
The identification of Lemma~\ref{lem:imag-root-mult-from-FA} is
conditional on the BL-2/4/6 comparison witnesses of the principal
locus (Hodge transport, BV residue transport, Maulik--Okounkov
stable-envelope transport) together with the oriented Hall/vanishing
cycle package. Within that locus one may ask for the elliptic
Feigin--Odesskii kernel on the Hall side to match the singular-theta
kernel after restriction to $E$. That match is the proof obligation.
Off the principal locus, or without the Hall package, there is no
global isomorphism and no reconstruction of $\Delta_{5}$ from hCS.
\end{heal}

\section{Non-perturbative completion}
\label{sec:non-perturbative-completion}

The perturbative $E_{3}$-FA $\Obs^{q,\hCS}(K3\times E,\frakg)$ does not
produce $\Delta_{5}$ by itself. It can at most supply a formal character
seed compatible with $\Delta_{5}^{-1}$ after the hCS--Hall/Borcherds
comparison is constructed.

\begin{remark}[Non-perturbative content of $\Delta_{5}$]
\label{rem:non-perturbative-delta5}
\ClaimStatusDerivation{}\\
The Borcherds product representation
\eqref{eq:borcherds-product-delta5} is a closed-form expression valid
on the full Siegel upper half-plane $\mathbb{H}_{2}$, including the
deep interior where $\Im Z$ is finite. The convergence of the
Borcherds product on $\mathbb{H}_{2}$ uses the modular invariance of
the singular-theta integral, which is non-perturbative in any natural
expansion parameter. By contrast, the perturbative $\hbar$-series of
$\Obs^{q,\hCS}$ is asymptotic in $\hbar$, conditional in $q$ to the
Klingen chamber. The completion from any compatible formal series to
the closed Borcherds product is therefore non-perturbative, supplied
externally either by the Borcherds singular-theta lift
\cite{Borcherds1998} or---equivalently on $K3\times E$---by the
Donaldson--Thomas generating function
$Z^{K3\times E}_{\mathrm{DT}}=C\cdot\Delta_{5}^{-2}$ of Oberdieck--%
Pixton \cite{OberdieckPixton2019}. The latter is a Donaldson--Thomas
quantity, not perturbative in any holomorphic Chern--Simons coupling,
and not visible in any finite Feynman-diagram expansion.
\end{remark}

\begin{proposition}[Three external completions compatible with the seed]
\label{prop:three-completions-perturbative-seed}
\ClaimStatusConditional{}\\
If the perturbative seed
$[\Obs^{q,\hCS}(K3\times E,\frakg)]_{\hbar\text{-formal}}$ is connected
to the Hall/Borcherds trace by the principal-locus witnesses, then it is
compatible with three external non-perturbative completions to
$\Delta_{5}(Z)$:
\begin{enumerate}
\item[\textup{(a)}] Borcherds singular-theta lift of the weak Jacobi
seed $\phi_{0,1}$ (\cite{Borcherds1998} Theorem 13.3), which gives the
closed Borcherds product;
\item[\textup{(b)}] Gritsenko additive lift of the Jacobi--cusp seed
$\eta^{9}\vartheta_{1}\in J_{9/2,1/2}^{\mathrm{cusp}}$
(\cite{GritsenkoNikulin1995}, \cite{Lorgat2020}), which converts the
half-integer-weight elliptic seed into the closed weight-$5$ Siegel
modular form $\Delta_{5}=\Grit(\eta^{9}\vartheta_{1})$;
\item[\textup{(c)}] Oberdieck--Pixton DT completion
\cite{OberdieckPixton2019}, which identifies
$Z^{K3\times E}_{\mathrm{DT}}=C\cdot\Delta_{5}^{-2}$ at the level of
DT moduli generating functions; the $\Delta_{5}^{-2}$ is the
non-perturbative DT-side counterpart, not a finite hCS Feynman output.
\end{enumerate}
The three completions agree on the principal locus by Borcherds'
multiplicative-additive comparison
\cite{Borcherds1998}, by the Gritsenko--Nikulin identity
$\Phi_{10}=\Delta_{5}^{2}$, and by Oberdieck--Pixton's matching of the
DT side to $\Phi_{10}^{-1}$ \cite{OberdieckPixton2019}.
\end{proposition}

\begin{proof}
For (a)--(b): the multiplicative-and-additive equivalence of the
Borcherds and Gritsenko lifts on the weight-$5$ Siegel modular line
$M_{5}(\Sp_{4}(\Z),v_{\Delta_{5}})=\C\cdot\Delta_{5}$ is
\cite{GritsenkoNikulin1995} Theorem 2.1 (their Theorem 2.1 establishes
$\Grit(\eta^{9}\vartheta_{1})=\BorLift(\phi_{0,1})$). For (c):
\cite{OberdieckPixton2019} Theorem 1 identifies the reduced DT generating
function with $1/\Phi_{10}=\Delta_{5}^{-2}$. Each completion supplies
the unique Siegel modular form of weight $5$ on $\Sp_{4}(\Z)$ with the
required multiplicities, hence the three external completions agree.
\end{proof}

\section{Status, scope, and remaining obstructions}
\label{sec:status-scope-remaining}

\textsc{Result.} On the principal $K3\times E$ locus of manuscript
Theorem~\texttt{thm:g-delta5-is-sp-k3}, the perturbative $\Ethree$-FA
$\Obs^{q,\hCS}(K3\times E,\frakg)$ does not reconstruct $\Delta_{5}$.
The stable typed statement is conditional compatibility:
\begin{enumerate}
\item[(R1)] (Lemma~\ref{lem:imag-root-mult-from-FA}) FA characters
become imaginary-root multiplicities only after the oriented
hCS--Hall/Borcherds comparison and vanishing-cycle normalisation are
supplied;
\item[(R2)] (Theorem~\ref{thm:perturbative-match-delta5}) the formal
$q$-series of FA characters has $\Delta_{5}(2Z)^{-1}|_{E}$ as its
conditional reciprocal-denominator target in $\Q[\![q,\hbar]\!]$;
\item[(R3)] (Theorem~\ref{thm:borcherds-weyl-rigidity-reconstruction})
Borcherds--Weyl rigidity gives the unique automorphic completion only
after the multiplicities are supplied externally.
\end{enumerate}

\textsc{Scope.} The compatibility is conditional on:
\begin{enumerate}
\item[(S1)] BL-2 (Hodge transport), BL-4 (BV residue transport), BL-6
(Maulik--Okounkov stable-envelope transport) on the principal
$K3\times E$ locus;
\item[(S2)] vanishing cubic Casimir $d^{abc}_{\frakg}=0$ for the gauge
dg-Lie input;
\item[(S3)] perturbative-and-formal scope in $\hbar$, with the
non-perturbative completion supplied externally by Borcherds 1998 /
Gritsenko 1999 / Oberdieck--Pixton 2019.
\end{enumerate}

\textsc{Open obstructions.}
\begin{enumerate}
\item[(O1)] The all-orders convergence of the perturbative
$\hbar$-series on the compact $K3\times E$ base is open (compact $6$d
analogue of \cite{CostelloYagi2018}); only the formal
$\hbar$-expansion is established.
\item[(O2)] The compatibility does not extend to the
$N\in\{2,3,4,6\}$ CHL siblings $\Phi_{N}$ without supplying the
order-$N$ equivariant oriented hCS-to-Hall comparison and the
corresponding $g_{N}$-equivariant specialisation cycle
$\Sigma_{2}^{(N)}$ (manuscript
\texttt{cor:sibling-BKMs-from-one-phiFA}, Conjecture
\texttt{conj:g-delta5-is-sp-k3-general-extension}).
\item[(O3)] The non-perturbative completion of
Proposition~\ref{prop:three-completions-perturbative-seed} is
classical---it does not produce $\Delta_{5}$ directly from
$\Obs^{q,\hCS}$ but invokes external automorphic input. The
Costello--Li-style direct identification of the Borcherds product as
the closed-form one-loop output, paralleling the Costello--Yagi $5$d
hCS-to-Yangian VOA all-orders theorem, is open in $6$d.
\end{enumerate}

The four $\kappa_{\bullet}$ values $\{0,3,5,24\}$ at $K3\times E$ are
preserved without conflation: $\kcat(K3\times E)=0$ (K\"unneth total
space, manuscript), $\kch^{\Heis}=3$ (Heisenberg--Cartan rank of
$\cA_{E}^{\hCS}$ on $E$, three Cartan directions), $\kBKM(\Gdelta)=5$
(weight of $\Delta_{5}$, Borcherds singular-theta), and $\kfib=24$
(Mukai-lattice rank of the K3 fibre).

\section*{Acknowledgements}

This note records the blocker-6 comparison against the manuscript
chapters \texttt{chapters/theory/cy\_to\_chiral.tex} and
\texttt{chapters/examples/k3e\_bkm\_chapter.tex}; the cited primary
sources are Gritsenko--Nikulin 1995, Borcherds 1998,
Oberdieck--Pixton 2019, Lorgat 2020.

\begin{thebibliography}{99}

\bibitem{Borcherds1998}
R.~E.~Borcherds,
\emph{Automorphic forms with singularities on Grassmannians},
Inventiones Mathematicae \textbf{132} (1998), 491--562.

\bibitem{Bruinier2002}
J.~H.~Bruinier,
\emph{Borcherds Products on $O(2,l)$ and Chern Classes of Heegner
Divisors},
Lecture Notes in Mathematics \textbf{1780}, Springer, 2002.

\bibitem{CostelloGwilliam2017}
K.~Costello, O.~Gwilliam,
\emph{Factorization Algebras in Quantum Field Theory, Vol.~1--2},
Cambridge University Press, 2017--2021.

\bibitem{CostelloYagi2018}
K.~Costello, M.~Yagi,
\emph{Unification of integrability in supersymmetric gauge theories},
Advances in Theoretical and Mathematical Physics \textbf{24} (2020),
1931--2041.

\bibitem{GritsenkoNikulin1995}
V.~A.~Gritsenko, V.~V.~Nikulin,
\emph{Igusa modular forms and ``the simplest'' Lorentzian Kac--Moody
algebras},
Mat.\ Sb.\ \textbf{187} (1996), 1601--1641; arXiv:alg-geom/9506016.

\bibitem{Kudla1997}
S.~Kudla,
\emph{Algebraic cycles on Shimura varieties of orthogonal type},
Duke Mathematical Journal \textbf{86} (1997), 39--78.

\bibitem{Lorgat2020}
R.~Lorgat,
\emph{Automorphic corrections and the Igusa cusp form $\Delta_{5}$},
preprint, 2020; arXiv:2003.~(automorphic corrections).

\bibitem{OberdieckPixton2019}
G.~Oberdieck, R.~Pixton,
\emph{Holomorphic anomaly equations and the Igusa cusp form conjecture},
Inventiones Mathematicae \textbf{213} (2018), 507--587.

\bibitem{SchiffmannVasserot2013}
O.~Schiffmann, E.~Vasserot,
\emph{The elliptic Hall algebra and the K-theory of the Hilbert scheme
of $\mathbb{A}^{2}$},
Duke Mathematical Journal \textbf{162} (2013), 279--366.

\end{thebibliography}

\end{document}
